% 2026-11-11 / v32
\section{独立中心三次変数の対称性によるArtin障害の消滅}
\label{csl:sec:symmetry}

ここでは論文Bの混合中心対において、$v$が独立中心変数であることを使う。
前節の五変数三次の体上の隣接sector計算を、\emph{正しい導来自然変換複体が構成された場合}の全Artin持上げ定理へ強化する。
これにより次数一の$2\times5$障害行列を逐次計算する必要がなくなるが、
導来bar複体の正当性やkernel--functor比較(KF)そのものは証明しない。

\subsection{体上の障害sectorへの補助対称性}
$k$を三乗根$\zeta$を含む標数零の体とし、
\[
 F_\lambda=(x^3+y^3+z^3)/3+u^3+v^3+\lambda uxy,
 \qquad \lambda^3\ne-3
\]
とする。$\Gamma=\mu_3^{(v)}$を$v\mapsto\zeta v$、
$x,y,z,u\mapsto x,y,z,u$で作用させる。
混合量子化の二次関係は$v$を含まず、指定中心元の$v$への依存は$v^3$のみである。
従って$\Gamma$は量子パラメータ$h$を固定し、非可換中心対へ作用する。

\begin{proposition}[障害sectorの非自明な$\Gamma$-指標]
\label{csl:prop:characters}
$\mathcal T_0$を$F_\lambda$のgraded MF圏とする。
$\Gamma$の作用について
\[
 H^0\RHom(\Delta,\Delta(1))
 =\langle x,y,z,u\rangle\oplus\langle v\rangle,
 \qquad
 \left(H^1\RHom(\Delta,\Delta(1))\right)^\Gamma=0.
 \tag{CSL1}\label{csl:eq:characters}
\]
最初の直和の左項は自明表現で、右項は非自明な一次元表現である。
後者の$H^1$は二つの非恒等scalar-$\mu_3$ sectorの和であり、
各sectorは$\Gamma$の非自明な一次元表現となる。
\end{proposition}
\begin{proof}
前節の命題\ref{aro:prop:sector-neighbors}から
$H^0\simeq\operatorname{Jac}(F_\lambda)_1$、
$H^1\simeq k^2$である。$H^0$の表現は変数の作用から従う。
後者の各非恒等sectorは原点のみを固定し、
normal空間は五つの座標方向全部である。
Ballard--Favero--Katzarkovのsector formula
\cite[Theorem~5.39, proof]{AROBFK}では、
normal方向のtop exterior power（双対規約ではその双対）の
行列式characterが現れる。$\Gamma$はnormalの$v$方向だけを
$\zeta$倍するので、その行列式characterは$\zeta^{\pm1}$であり
自明でない。$\Gamma$はscalar-$\mu_3$と可換で各sectorを保つから、
二つのsectorの和にも不変元はない。
この判定にはcharacterの双対規約を選ぶ必要はない。
\end{proof}

\subsection{有限Artin環上の一般持上げ補題}
$R$を剰余体$k$を持つ局所Artin $k$-代数とし、$\Gamma$は$R$を固定する。
以下の$R$-線形cochain複体$C_R$は各項が$R$-平坦で、
すべてのArtin商への縮小と\emph{鎖レベルで}両立すると仮定する。
実際のdg圏の導来自然変換を正しく表すという仮定は、この条件と別に必要である。

\begin{lemma}[対称性と一つの恒等的閉元による全Artin自由性]
\label{csl:lem:equivariant-lift}
$C_R$を上記の$\Gamma$-同変な複体の族とし、特殊ファイバー$C_k$について
\[
 H^{-1}(C_k)=0,\quad H^0(C_k)=V_0\oplus kv,
 \quad V_0=H^0(C_k)^\Gamma\simeq k^4,
 \quad H^1(C_k)^\Gamma=0
\]
を仮定する。$v$の類はすべての$R$で互換な閉元$v_R\in Z^0(C_R)$へ
持ち上がるとする。このとき、
\[
 H^{-1}(C_R)=0,\qquad H^0(C_R)\simeq R^5,
 \qquad H^0(C_R)\otimes_R R/I\simeq H^0(C_{R/I})
 \tag{CSL2}\label{csl:eq:artin-free}
\]
が任意のArtin商$R/I$に対して成り立つ。
\end{lemma}
\begin{proof}
$R$の長さについて帰納し、長さ一の小拡大
$0\to I\to R'\to R\to0$を取る。ただし$\mathfrak m_{R'}I=0$。
項別平坦性から
\[
 0\to I\otimes_k C_k\to C_{R'}\to C_R\to0
\]
は複体の短完全列となる。$H^{-1}(C_k)=0$より、
帰納法で$H^{-1}(C_{R'})=0$を得る。

$V_0$の基底の不変閉代表を$C_R$に持ち上げてあるとする。
$\Gamma$は有限群で$\operatorname{char}k=0$だから、
不変部分を取る操作は完全であり、$C_{R'}^0$に不変な持上げを選べる。
その微分は$I\otimes C_k^1$に属する不変なcocycleである。
仮定$H^1(C_k)^\Gamma=0$より、
$I\otimes C_k^0$の不変な元の微分として消せる。
従って四基底は閉元のまま持ち上がる。
$v_R$には初めから互換な閉元がある。
よって$H^0(C_{R'})\to H^0(C_R)$は全射。
長完全列と$H^{-1}(C_R)=0$により
\[
 0\to I\otimes_k H^0(C_k)\to H^0(C_{R'})
 \to H^0(C_R)\to0
\]
が成り立つ。従って$H^0(C_{R'})$の$k$-長さは
$5\operatorname{length}(R')$であり、剰余体の五基底の持上げから
$R'^5\twoheadrightarrow H^0(C_{R'})$が得られる。
Nakayamaと長さの一致により同型である。
小拡大について繰り返せば任意の$R$とすべての商で結論する。
\end{proof}

\begin{theorem}[混合族の低rotation底変更に関する条件付き対称性定理]
\label{csl:thm:conditional-bc}
$F_\lambda$を特殊ファイバーとする$\Gamma$-同変な
Artin変形のgraded MFあるいは非可換MFについて、
$r=1,2$の\emph{実際の導来自然変換}を計算する$R$-線形cochain複体$C_{r,R}$が
次の条件を満たすとする：
(i) 項別$R$-平坦、任意の商への剰余が鎖レベルで一致する；
(ii) $\Gamma$同変で、特殊ファイバーでは既知のsector同定と作用に一致する；
(iii) 中心変数$v$による積が$C_{1,R}$の互換な閉元を与える。
このとき、任意の$R$に対し
\[
 H^0(C_{1,R})\simeq R^5,\qquad
 H^0(C_{2,R})\simeq R^{10},
 \tag{CSL3}\label{csl:eq:bc}
\]
かつ両者のすべてのArtin底変更が同型である。
特に論文Bの乗法的characteristic写像$\Theta_r$の
自然な特殊化が既知の多項式倍との同定であるなら、
$\Theta_1,\Theta_2$は局所環上の単元行列式判定によって同型となり、
条件(BC)が満たされる。
\end{theorem}
\begin{proof}
$r=1$には命題\ref{csl:prop:characters}と
補題\ref{csl:lem:equivariant-lift}を使う。
$r=2$では前節の\ref{aro:prop:sector-neighbors}から
$H^{-1}(C_{2,k})=H^1(C_{2,k})=0$、
$H^0(C_{2,k})\simeq k^{10}$であり、
補題\ref{aro:lem:neighbor-vanishing}の小拡大版を適用する。
積と$\Theta$の条件付き同定は、標的とsourceの自由性および特殊ファイバーでの同型から、$\det(\Theta_r)$が単元となることによる。
\end{proof}

\begin{remark}[解決したこと・残ること]\label{csl:rem:scope}
この定理は\emph{仮定した正しい複体の内部での}全次数の障害消滅であり、
任意の量子化や一般混合パラメータで条件(i)--(iii)が成立するとは主張しない。
特に非正則Artin上の導来二側bar、相対cofibrant模型、
$\Theta$の実際の導来自然変換への写像、kernel全忠実性(KF)の検証は
別に残る。これらを体上の次元公式や項別平坦性から推論しない。
一方、$\Gamma$が保存される部分族については、
以前の$2\times5$障害行列の高次係数を個別に計算する優先度は下がり、
正しいArtin互換のenhancementの構成が最優先となる。
\end{remark}
